%
% 53-verallgemeinert.tex -- Das verallgemeinerte Argumentprinzip
%
% (c) 2026 Prof Dr Andreas Müller
%
\subsection{Das verallgemeinerte Argumentprinzip
\label{buch:analytisch:argument:subsection:verallgemeinert}}
Das Argumentprinzip von Satz~\ref{buch:analytisch:argument:satz:argument}
hat eine unmittelbare Verallgemeinerung.

\begin{satz}[Verallgemeinertes Argumentprinzip]
\index{Argumentprinzip!verallgemeinert}%
Sei $f$ eine meromorphe Funktion und $\gamma$ ein Weg, der keine Nullstellen
oder Pole von $f$ enthält und sei $g$ eine ganze Funktion.
Dann ist
\[
\oint_\gamma \frac{f'(z)}{f(z)} g(z)\,dz
=
\sum_{\text{$a$ Nullstelle von $f$}} \operatorname{ind}_\gamma(a) g(a)
-
\sum_{\text{$b$ Pol von $f$}} \operatorname{ind}_\gamma(b) g(a)
\]
Wobei die Nullstellen $a$ und Pole $b$ von $f$ mit Ihrer Vielfachheit
gezählt werden sollen.
\end{satz}

